% !TEX root = ../../main.tex

\section{Probability Exercises}

\subsection{Basic Probability Concepts}

\begin{enumerate}
    \item Define probability and explain the difference between classical, empirical, and subjective probability.
    
    \item A fair six-sided die is rolled. What is the probability of rolling an even number?
    
    \item If two events A and B are mutually exclusive, what is P(A or B)?
    
    \item Explain the addition rule for probabilities. Provide an example.
    
    \item What is conditional probability? Give the formula and an example.
\end{enumerate}

\subsection{Bayes' Theorem and Conditional Probability}

\begin{enumerate}
    \item State Bayes' Theorem and explain its significance.
    
    \item In a population, 1\% of people have a certain disease. A test for the disease is 99\% accurate for both positive and negative results. If a person tests positive, what is the probability they actually have the disease?
    
    \item Two cards are drawn without replacement from a standard deck of 52 cards. What is the probability that both are aces?
    
    \item Explain the multiplication rule for independent events. Provide an example.
    
    \item If P(A|B) = 0.6, P(B) = 0.4, and P(A and B) = 0.24, find P(B|A).
\end{enumerate}


\subsection{Random Variables and Distributions}

\begin{enumerate}
    \item Define a random variable. Distinguish between discrete and continuous random variables.
    
    \item For a binomial distribution with n=10 trials and p=0.5, calculate the probability of exactly 5 successes.
    
    \item What is the mean and variance of a Poisson distribution with parameter $\lambda=3$?
    
    \item Explain the normal distribution. What are its key properties?
    
    \item A random variable X follows a normal distribution with mean 50 and standard deviation 10. Find P(X > 60).
\end{enumerate}

\subsection{Review Questions for Tests}

\begin{enumerate}
    \item Solve the following: In a box of 20 balls, 5 are red, 10 are blue, and 5 are green. Two balls are drawn at random without replacement. Find the probability that both are blue.
    
    \item Use Bayes' Theorem to solve: A factory produces items, 90\% of which are good. A quality control test correctly identifies good items 95\% of the time and bad items 80\% of the time. If an item passes the test, what is the probability it is actually good?
    
    \item For a geometric distribution with p=0.2, find the probability that the first success occurs on the 3rd trial.
    
    \item Calculate the expected value and variance for a binomial random variable with n=8, p=0.3.
    
    \item In a normal distribution with $\mu=100$ and $\sigma=15$, find the z-score for X=120 and interpret it.
\end{enumerate}

\section{Solutions to Exercises}

\subsection{Basic Probability Concepts}

\begin{enumerate}
    \item \textbf{Definition of Probability and Its Interpretations:}
    \begin{itemize}
        \item \textbf{Classical Probability:} If there are $n$ equally likely outcomes and $m$ favorable outcomes, $P(A) = \frac{m}{n}$. Assumes all outcomes are equally likely. Example: rolling a fair die.
        \item \textbf{Empirical (Frequentist) Probability:} Based on observed frequencies in repeated experiments. $P(A) = \lim_{N \to \infty} \frac{\text{number of times A occurs}}{N}$. Requires data collection.
        \item \textbf{Subjective Probability:} A person's belief or degree of confidence that an event will occur. Differs between individuals based on information and experience.
    \end{itemize}
    
    \item \textbf{Rolling an Even Number on a Fair Die:}
    Even numbers: $\{2, 4, 6\}$. Total outcomes: 6.
    \begin{equation}
        P(\text{even}) = \frac{3}{6} = \frac{1}{2} = 0.5
    \end{equation}
    
    \item \textbf{Mutually Exclusive Events:}
    If events A and B are mutually exclusive, they cannot occur simultaneously. Therefore:
    \begin{equation}
        P(A \text{ or } B) = P(A) + P(B)
    \end{equation}
    (No overlap to subtract.)
    
    \item \textbf{Addition Rule for Probabilities:}
    \begin{equation}
        P(A \text{ or } B) = P(A) + P(B) - P(A \text{ and } B)
    \end{equation}
    \textit{Example:} Rolling a die and getting a number greater than 2 OR even.
    \begin{itemize}
        \item $P(\text{greater than 2}) = \frac{4}{6}$ (outcomes: 3,4,5,6)
        \item $P(\text{even}) = \frac{3}{6}$ (outcomes: 2,4,6)
        \item $P(\text{both}) = \frac{2}{6}$ (outcomes: 4,6)
        \item $P(\text{greater than 2 or even}) = \frac{4}{6} + \frac{3}{6} - \frac{2}{6} = \frac{5}{6}$
    \end{itemize}
    
    \item \textbf{Conditional Probability:}
    Conditional probability is the probability of event A occurring given that event B has already occurred:
    \begin{equation}
        P(A|B) = \frac{P(A \text{ and } B)}{P(B)}
    \end{equation}
    \textit{Example:} If P(rain) = 0.3 and P(traffic jam | rain) = 0.8, the probability of a traffic jam given rain is 0.8.
\end{enumerate}

\subsection{Bayes' Theorem and Conditional Probability}

\begin{enumerate}
    \item \textbf{Bayes' Theorem:}
    \begin{equation}
        P(A|B) = \frac{P(B|A) \cdot P(A)}{P(B)}
    \end{equation}
    \textit{Significance:} Allows us to update beliefs based on new evidence. Central to probability, statistics, machine learning, and decision-making.
    
    \item \textbf{Disease Test Problem:}
    \begin{itemize}
        \item Let D = disease present, + = positive test
        \item P(D) = 0.01, P(D$^c$) = 0.99
        \item P(+|D) = 0.99, P(+|D$^c$) = 0.01
        \item Find: P(D|+)
    \end{itemize}
    Using Bayes:
    \begin{equation}
        P(D|+) = \frac{P(+|D) \cdot P(D)}{P(+)} = \frac{0.99 \times 0.01}{0.99 \times 0.01 + 0.01 \times 0.99} = \frac{0.0099}{0.0198} \approx 0.50
    \end{equation}
    \textbf{Result:} Despite a 99\% accurate test, only ~50\% probability of having the disease if tested positive. This illustrates how base rates matter.
    
    \item \textbf{Two Aces Without Replacement:}
    \begin{itemize}
        \item First card is an ace: $P(\text{ace}_1) = \frac{4}{52}$
        \item Second card is an ace given first was an ace: $P(\text{ace}_2 | \text{ace}_1) = \frac{3}{51}$
    \end{itemize}
    \begin{equation}
        P(\text{both aces}) = \frac{4}{52} \times \frac{3}{51} = \frac{12}{2652} = \frac{1}{221} \approx 0.0045
    \end{equation}
    
    \item \textbf{Multiplication Rule for Independent Events:}
    If events A and B are independent:
    \begin{equation}
        P(A \text{ and } B) = P(A) \times P(B)
    \end{equation}
    \textit{Example:} Rolling two dice. P(first die = 3 AND second die = 4) = $\frac{1}{6} \times \frac{1}{6} = \frac{1}{36}$.
    
    \item \textbf{Finding P(B|A):}
    Using the definition of conditional probability:
    \begin{equation}
        P(A|B) = \frac{P(A \text{ and } B)}{P(B)}
    \end{equation}
    Rearranging:
    \begin{equation}
        P(A \text{ and } B) = P(A|B) \times P(B) = 0.6 \times 0.4 = 0.24
    \end{equation}
    (This matches the given information.) Now find P(B|A):
    \begin{equation}
        P(B|A) = \frac{P(A \text{ and } B)}{P(A)} = \frac{0.24}{P(A)}
    \end{equation}
    We need P(A). Using P(A and B) = 0.24 and P(A|B) = 0.6, we have P(A) = 0.6 × 0.4 / 0.24 × ... wait, let me recalculate.
    We have: $P(A|B) \times P(B) = P(A \text{ and } B)$, so $0.6 \times 0.4 = 0.24$ 
    We need another equation. Using $P(B|A) = \frac{P(A \text{ and } B)}{P(A)}$, but we don't have P(A) directly.
    Actually, we can derive: $P(A) = \frac{P(A \text{ and } B)}{P(B|A)}$. But we're looking for P(B|A), so:
    \begin{equation}
        P(B|A) = \frac{P(A \text{ and } B)}{P(A)}
    \end{equation}
    Without additional information to determine P(A), we need more constraints. However, if we assume we need to find P(B|A) in terms of given information: Since events might be related, we'd need P(A) to complete this calculation. Assuming independence or additional context, a reasonable answer depends on the problem setup.
\end{enumerate}

\subsection{Random Variables and Distributions}

\begin{enumerate}
    \item \textbf{Random Variables:}
    A random variable is a function that assigns a numerical value to each outcome in a sample space.
    \begin{itemize}
        \item \textbf{Discrete:} Takes on countable values (e.g., 0, 1, 2, ...). Example: number of heads in 10 coin flips.
        \item \textbf{Continuous:} Takes on any value in an interval. Example: height of a person, temperature.
    \end{itemize}
    
    \item \textbf{Binomial Probability (n=10, p=0.5, exactly 5 successes):}
    \begin{equation}
        P(X = 5) = \binom{10}{5} (0.5)^5 (0.5)^5 = 252 \times (0.5)^{10} = \frac{252}{1024} \approx 0.246
    \end{equation}
    
    \item \textbf{Poisson Distribution with $\lambda = 3$:}
    \begin{equation}
        E[X] = \lambda = 3
    \end{equation}
    \begin{equation}
        \text{Var}(X) = \lambda = 3
    \end{equation}
    
    \item \textbf{Normal Distribution Properties:}
    \begin{itemize}
        \item Symmetric around the mean $\mu$
        \item Bell-shaped curve
        \item Defined by two parameters: mean $\mu$ and standard deviation $\sigma$
        \item About 68\% of data within $\pm 1\sigma$, 95\% within $\pm 2\sigma$, 99.7\% within $\pm 3\sigma$
        \item Many natural phenomena approximately follow normal distribution
    \end{itemize}
    
    \item \textbf{Normal Distribution: P(X > 60) where $\mu=50$, $\sigma=10$:}
    Standardize to z-score:
    \begin{equation}
        Z = \frac{X - \mu}{\sigma} = \frac{60 - 50}{10} = 1
    \end{equation}
    From standard normal table: $P(Z > 1) \approx 0.1587$ or 15.87\%
\end{enumerate}

\subsection{Review Questions for Tests - Solutions}

\begin{enumerate}
    \item \textbf{Two Blue Balls Without Replacement:}
    \begin{itemize}
        \item Total balls: 20, Blue balls: 10
        \item First blue: $P(\text{blue}_1) = \frac{10}{20}$
        \item Second blue given first blue: $P(\text{blue}_2 | \text{blue}_1) = \frac{9}{19}$
    \end{itemize}
    \begin{equation}
        P(\text{both blue}) = \frac{10}{20} \times \frac{9}{19} = \frac{90}{380} = \frac{9}{38} \approx 0.237
    \end{equation}
    
    \item \textbf{Quality Control - Bayes' Theorem:}
    \begin{itemize}
        \item Let G = good item, + = passes test
        \item P(G) = 0.90, P(G$^c$) = 0.10
        \item P(+|G) = 0.95 (95\% correct for good items)
        \item P(+|G$^c$) = 0.20 (80\% correct for bad items means 20\% false positive)
    \end{itemize}
    Total probability of passing:
    \begin{equation}
        P(+) = P(+|G) \cdot P(G) + P(+|G^c) \cdot P(G^c) = 0.95 \times 0.90 + 0.20 \times 0.10 = 0.855 + 0.02 = 0.875
    \end{equation}
    Bayes' Theorem:
    \begin{equation}
        P(G|+) = \frac{P(+|G) \cdot P(G)}{P(+)} = \frac{0.95 \times 0.90}{0.875} = \frac{0.855}{0.875} \approx 0.977
    \end{equation}
    
    \item \textbf{Geometric Distribution (p=0.2, first success on trial 3):}
    \begin{equation}
        P(X = 3) = (1-p)^{3-1} \cdot p = (0.8)^2 \times 0.2 = 0.64 \times 0.2 = 0.128
    \end{equation}
    
    \item \textbf{Binomial Expected Value and Variance (n=8, p=0.3):}
    \begin{equation}
        E[X] = np = 8 \times 0.3 = 2.4
    \end{equation}
    \begin{equation}
        \text{Var}(X) = np(1-p) = 8 \times 0.3 \times 0.7 = 1.68
    \end{equation}
    
    \item \textbf{Normal Distribution z-score ($\mu=100$, $\sigma=15$, X=120):}
    \begin{equation}
        Z = \frac{X - \mu}{\sigma} = \frac{120 - 100}{15} = \frac{20}{15} = 1.33
    \end{equation}
    \textbf{Interpretation:} The value 120 is 1.33 standard deviations above the mean. This is approximately at the 91.8th percentile (about 92\% of data falls below this value). It represents an above-average outcome but not extremely rare.
\end{enumerate}